\section{Related Work}
\label{sec:related}

This paper is closest to work on text watermarking, LLM fingerprinting,
publicly auditable watermark protocols, provenance standards, proof systems,
model ownership, and extraction diagnostics.  We organize the literature by
the substrate that carries evidence, because substrate choice determines which
verification properties can be claimed.

\paragraph{Text watermarking and public final text.}
Watermarking methods bias or test generated text so that a detector can
distinguish watermarked samples from unwatermarked samples
\cite{kirchenbauer2023watermark,zhao2023provable,hu2024unbiased,kuditipudi2024robust,christ2023undetectable,lee2024who,chang2024postmark,huang2025waterpool,dabiriaghdam2025simmark,wouters2024optimizing,liu2024adaptive,zhu2024duwak,lau2024waterfall,jiang2025stealthink}.
Production-scale systems such as SynthID-Text show that text watermarking can
be deployed at large scale with a detector-oriented claim
\cite{dathathri2024scalable}.  This family treats final text as the main
evidence object, so it is directly relevant to portability and
non-cooperation.  It also exposes the central substrate risk studied here: if
a public final-text predicate is observable, an attacker may search, edit,
paraphrase, or optimize against it.  The present manuscript does not claim
watermark success, public text-only verification success, or ownership proof;
it studies why the tested public final-text predicates recover no codeword
blocks and are spoofable.

\paragraph{Publicly auditable and protocol-backed watermarks.}
Recent work moves beyond private-key text detectors by designing publicly
detectable or publicly verifiable watermark protocols
\cite{fairoze2023publicly,isler2023puppy,duan2025pvmark,luan2026vow}.
These systems strengthen the verification interface by adding public
detectability, commitments, zero-knowledge style checks, oblivious detection,
or protocol participation.  In the VSG taxonomy, that is a substrate shift:
the evidence is no longer ordinary final text alone, but final text plus a
stronger protocol or proof object.  This paper is not a replacement for those
constructions.  It explains why a provider-side trace-bound diagnostic should
not be promoted into a final-text-only claim unless the stronger substrate is
actually available.

\paragraph{Metadata, credentials, and proof substrates.}
Content provenance standards such as C2PA Content Credentials bind assertions,
claims, and signatures into manifests associated with digital assets
\cite{c2pa2025spec}.  Proof-oriented LLM work such as zkLLM and ZKPROV studies
zero-knowledge verification for model computation or dataset provenance
\cite{sun2024zkllm,namazi2025zkprov}.  These approaches make evidence stronger
by moving it to a metadata, protocol, or proof substrate.  Cryptographic
provenance is outside the evidence scope of this manuscript; these systems are
contrast cases showing how much stronger the evidence object becomes when it
does not have to survive as ordinary natural text.

\paragraph{Learnability, removal, and spoofing.}
Removal and residual-signal studies emphasize that provenance detectors should
not be conflated with deterministic codeword recovery
\cite{chen2025demark,sander2024radioactive}.  Impossibility and attack work
also shows that strong claims can fail when the attacker can transform,
imitate, or learn watermark signals
\cite{zhang2023watermarkssand,gu2023learnability,an2025ditto}.  Our attack
ladder has the same cautionary role, but with a narrower target: it does not
claim a universal impossibility theorem over all public predicates.  It shows
that the tested public final-text predicates are in a weak position because
they recover zero codeword blocks while still offering source-mismatch accept
regions that can be searched, edited, learned, and directly optimized.

\paragraph{Active fingerprints and model-side behavior.}
Active fingerprints query a model with selected prompts and verify associated
responses or behaviors
\cite{xu2024instructional,zeng2024huref,russinovich2024chainhash,pasquini2025llmmap,yang2024fingerprint,yamabe2025mergeprint,xu2025ctcc,xu2025evertracer,bai2025esf,hu2025resf,nasery2025scalable,tsai2025rofl}.
These methods are closer to model ownership disputes than generic text
watermarks because they can involve owner-chosen prompts or behaviors.  Their
verification substrate, however, is often a response pattern, classifier score,
or query transcript rather than ordinary copied final text alone.  Our work
does not claim superiority over these methods.  It asks a narrower diagnostic
question: when evidence is visible in provider-side token events but not in
public final text, what can and cannot be verified?

\paragraph{Model-state and training-time ownership evidence.}
Classical DNN watermarking embeds signatures into parameters, activations, or
trigger responses
\cite{uchida2017embedding,zhang2018protecting,adi2018turning,rouhani2019deepsigns,jia2021entangled,szyller2021dawn,krauss2024clearstamp}.
These substrates can be strong when a verifier has white-box or cooperative
access, but they are less portable with copied text.  Scrubbing, backdoor
removal, model merging, and task-vector methods show that model-side evidence
can change under post-release transformation
\cite{zhang2025meraser,zeng2024beear,li2024backdoor,hubinger2024sleeper,wortsman2022model,ilharco2023task,yadav2023ties}.
The present paper does not claim robustness to arbitrary fine-tuning,
semantic rewriting, or probe-free distillation.

\paragraph{Passive provenance and extraction.}
Model extraction, memorization, and data-provenance studies infer exposure or
lineage without necessarily injecting a recoverable owner payload
\cite{carlini2024stealing,liang2024model,xie2024recall,puerto2025scaling,ravichander2025imprint}.
They motivate the non-cooperative dispute setting, but they usually answer a
different question from deterministic codeword recovery.  Our public
final-text experiments should therefore be read as a substrate diagnostic, not
as a claim that all passive provenance questions are impossible.

\paragraph{Positioning.}
The contribution here is neither a new public text-only verifier nor an
ownership-proof construction.  It is a claim-controlled map of which substrate
currently carries evidence.  Provider-side first-divergence traces show
controllability under event access; public final-text predicates in the
current artifacts fail to recover codewords and become source-mismatch
spoofing targets.  This framing complements prior methods by making explicit
when a method's evidence object travels with ordinary natural text and when it
requires a stronger cooperative, metadata, protocol, proof, or model-state
substrate.
